% Chunk: \S7.4 Lorentz Invariance
% Source: Weinberg QFT Vol. I, physical pp. 337--341 (printed pp. 314--318; begins and ends mid-page)
% Status: transcribed and source-reviewed; equations (7.4.1)--(7.4.26); no figures; no uncertainties

\subsection{Lorentz Invariance}

We are now going to show that the Lorentz invariance of the Lagrangian density
implies the Lorentz invariance of the $S$-matrix. Consider an infinitesimal
Lorentz transformation
\begin{equation}
\Lambda^\mu{}_\nu=\delta^\mu{}_\nu+\omega^\mu{}_\nu,
\tag{7.4.1}\label{eq:7.4.1}
\end{equation}
\begin{equation}
\omega_{\mu\nu}=-\omega_{\nu\mu}.
\tag{7.4.2}\label{eq:7.4.2}
\end{equation}
According to the analysis of the previous section, the invariance of the
action under such transformations tells us immediately that there are a set of
conserved `currents' $\mathcal{M}^{\rho\mu\nu}$:
\begin{equation}
\partial_\rho\mathcal{M}^{\rho\mu\nu}=0,
\tag{7.4.3}\label{eq:7.4.3}
\end{equation}
\begin{equation}
\mathcal{M}^{\rho\mu\nu}=-\mathcal{M}^{\rho\nu\mu},
\tag{7.4.4}\label{eq:7.4.4}
\end{equation}
one current for each independent component of $\omega_{\mu\nu}$. The integrals
of the time-components of these `currents' then provide us with a set of
time-independent tensors:
\begin{equation}
J^{\mu\nu}\equiv\int d^3x\,\mathcal{M}^{0\mu\nu},
\tag{7.4.5}\label{eq:7.4.5}
\end{equation}
\begin{equation}
\frac{d}{dt}J^{\mu\nu}=0.
\tag{7.4.6}\label{eq:7.4.6}
\end{equation}
The $J^{\mu\nu}$ will turn out to be the generators of the homogeneous Lorentz
group.

We would like to have an explicit formula for the tensor
$\mathcal{M}^{\rho\mu\nu}$, but Lorentz transformations act on the coordinates
and hence cannot leave the Lagrangian density invariant, so we cannot
immediately use the results of the previous section. However, translation
invariance allows us to formulate Lorentz invariance as a symmetry of the
Lagrangian density under a set of transformations on the fields and
field-derivatives alone. The fields undergo the matrix transformation
\begin{equation}
\delta\Psi^\ell=\frac{i}{2}\omega^{\mu\nu}
(\mathcal{J}_{\mu\nu})^\ell{}_m\Psi^m,
\tag{7.4.7}\label{eq:7.4.7}
\end{equation}
where $\mathcal{J}_{\mu\nu}$ are a set of matrices satisfying the algebra of the
homogeneous Lorentz group
\begin{equation}
[\mathcal{J}_{\mu\nu},\mathcal{J}_{\rho\sigma}]
=i\mathcal{J}_{\rho\nu}\eta_{\mu\sigma}
-i\mathcal{J}_{\sigma\nu}\eta_{\mu\rho}
-i\mathcal{J}_{\mu\rho}\eta_{\nu\sigma}
+i\mathcal{J}_{\mu\sigma}\eta_{\nu\rho}.
\tag{7.4.8}\label{eq:7.4.8}
\end{equation}
For example, for a scalar field $\phi$ we have $\delta\phi=0$, so
$\mathcal{J}_{\mu\nu}=0$, while for an irreducible field of type $(A,B)$ we have
\[
\mathcal{J}_{ij}=\epsilon_{ijk}(\mathcal{A}_k+\mathcal{B}_k),
\qquad
\mathcal{J}_{k0}=-i(\mathcal{A}_k-\mathcal{B}_k),
\]
where $\mathcal{A}$ and $\mathcal{B}$ are spin matrices for spin $A$ and $B$,
respectively. We specially note that for a covariant vector field,
$\delta V_\kappa=\omega_\kappa{}^\lambda V_\lambda$, so here
\[
(\mathcal{J}_{\mu\nu})_\kappa{}^\lambda
=-i\eta_{\mu\kappa}\delta_\nu{}^\lambda
+i\eta_{\nu\kappa}\delta_\mu{}^\lambda.
\]
The derivative of a field that transforms as in Eq.~\eqref{eq:7.4.7} transforms like
another such field, but with an extra vector index
\begin{equation}
\delta(\partial_\kappa\Psi^\ell)
=\frac{i}{2}\omega^{\mu\nu}(\mathcal{J}_{\mu\nu})^\ell{}_m
\partial_\kappa\Psi^m
+\omega_\kappa{}^\lambda\partial_\lambda\Psi^\ell.
\tag{7.4.9}\label{eq:7.4.9}
\end{equation}
The Lagrangian density is assumed to be invariant under the combined
transformations Eqs.~\eqref{eq:7.4.7} and \eqref{eq:7.4.9}, so
\begin{align*}
0={}&\frac{\partial\mathcal{L}}{\partial\Psi^\ell}
\frac{i}{2}\omega^{\mu\nu}(\mathcal{J}_{\mu\nu})^\ell{}_m\Psi^m
+\frac{\partial\mathcal{L}}{\partial(\partial_\kappa\Psi^\ell)}
\frac{i}{2}\omega^{\mu\nu}(\mathcal{J}_{\mu\nu})^\ell{}_m
\partial_\kappa\Psi^m\\
&+\frac{\partial\mathcal{L}}{\partial(\partial_\kappa\Psi^\ell)}
\omega_\kappa{}^\lambda\partial_\lambda\Psi^\ell.
\end{align*}
Setting the coefficient of $\omega^{\mu\nu}$ equal to zero gives
\begin{align*}
0={}&\frac{i}{2}\frac{\partial\mathcal{L}}{\partial\Psi^\ell}
(\mathcal{J}_{\mu\nu})^\ell{}_m\Psi^m
+\frac{i}{2}\frac{\partial\mathcal{L}}
{\partial(\partial_\kappa\Psi^\ell)}
(\mathcal{J}_{\mu\nu})^\ell{}_m\partial_\kappa\Psi^m\\
&+\frac12\frac{\partial\mathcal{L}}
{\partial(\partial_\kappa\Psi^\ell)}
(\eta_{\mu\kappa}\partial_\nu-\eta_{\nu\kappa}\partial_\mu)\Psi^\ell.
\end{align*}
Using the Euler--Lagrange equations~\eqref{eq:7.2.9}, and our formula~\eqref{eq:7.3.34} for the
energy--momentum tensor $T_{\mu\nu}$, we may write this as
\begin{equation}
0=\partial_\kappa\left[
\frac{i}{2}\frac{\partial\mathcal{L}}
{\partial(\partial_\kappa\Psi^\ell)}
(\mathcal{J}_{\mu\nu})^\ell{}_m\Psi^m\right]
-\frac12(T_{\mu\nu}-T_{\nu\mu}).
\tag{7.4.10}\label{eq:7.4.10}
\end{equation}
This immediately suggests the definition of a new energy--momentum tensor,
known as the \emph{Belinfante tensor}~\hyperref[ref:7.2]{[2]}:
\begin{align}
\Theta^{\mu\nu}=T^{\mu\nu}-\frac{i}{2}\partial_\kappa\Bigg[&
\frac{\partial\mathcal{L}}{\partial(\partial_\kappa\Psi^\ell)}
(\mathcal{J}^{\mu\nu})^\ell{}_m\Psi^m
-\frac{\partial\mathcal{L}}{\partial(\partial_\mu\Psi^\ell)}
(\mathcal{J}^{\kappa\nu})^\ell{}_m\Psi^m\notag\\
&-\frac{\partial\mathcal{L}}{\partial(\partial_\nu\Psi^\ell)}
(\mathcal{J}^{\kappa\mu})^\ell{}_m\Psi^m\Bigg].
\tag{7.4.11}\label{eq:7.4.11}
\end{align}
The quantity in square brackets is manifestly antisymmetric in $\mu$ and
$\kappa$, so $\Theta^{\mu\nu}$ satisfies the same conservation law as
$T^{\mu\nu}$,
\begin{equation}
\partial_\mu\Theta^{\mu\nu}=0.
\tag{7.4.12}\label{eq:7.4.12}
\end{equation}
For the same reason, when we set $\mu=0$ in Eq.~\eqref{eq:7.4.11} the index $\kappa$
runs over space components only, so the derivative term here drops out when we
integrate over all space
\begin{equation}
\int\Theta^{0\nu}d^3x=\int T^{0\nu}d^3x=P^\nu,
\tag{7.4.13}\label{eq:7.4.13}
\end{equation}
where $P^0\equiv H$. Thus $\Theta^{\mu\nu}$ can be regarded as \emph{the}
energy--momentum tensor, just as well as $T^{\mu\nu}$. However, Eq.~\eqref{eq:7.4.10}
tells us that, unlike $T^{\mu\nu}$, the Belinfante tensor $\Theta^{\mu\nu}$ is
not only conserved but also \emph{symmetric}:
\begin{equation}
\Theta^{\mu\nu}=\Theta^{\nu\mu}.
\tag{7.4.14}\label{eq:7.4.14}
\end{equation}
It is $\Theta^{\mu\nu}$ rather than $T^{\mu\nu}$ that acts as a source of the
gravitational field~\hyperref[ref:7.3]{[3]}. In consequence of the symmetry of $\Theta^{\mu\nu}$,
we may construct one more conserved tensor density:
\begin{equation}
\mathcal{M}^{\lambda\mu\nu}
\equiv x^\mu\Theta^{\lambda\nu}-x^\nu\Theta^{\lambda\mu}.
\tag{7.4.15}\label{eq:7.4.15}
\end{equation}
This is conserved, in the sense that
\begin{equation}
\partial_\lambda\mathcal{M}^{\lambda\mu\nu}
=\Theta^{\mu\nu}-\Theta^{\nu\mu}=0.
\tag{7.4.16}\label{eq:7.4.16}
\end{equation}
Thus Lorentz invariance allows us to define one more time-independent tensor
\begin{equation}
J^{\mu\nu}=\int\mathcal{M}^{0\mu\nu}d^3x
=\int d^3x\bigl(x^\mu\Theta^{0\nu}-x^\nu\Theta^{0\mu}\bigr).
\tag{7.4.17}\label{eq:7.4.17}
\end{equation}
The rotation generator $J_k=\epsilon_{ijk}J^{ij}/2$ is not only
time-independent, but also has no explicit time-dependence, so it commutes with
the Hamiltonian
\begin{equation}
[H,\mathbf{J}]=0.
\tag{7.4.18}\label{eq:7.4.18}
\end{equation}
Also, applying Eq.~\eqref{eq:7.3.28} to the function $\Theta^{0\nu}$, we have
\begin{align*}
[P_j,J_i]
&=\frac12\epsilon_{k\ell i}[P_j,J^{k\ell}]\\
&=\frac{i}{2}\epsilon_{k\ell i}\int d^3x\left(
x^k\frac{\partial}{\partial x^j}\Theta^{0\ell}
-x^\ell\frac{\partial}{\partial x^j}\Theta^{0k}\right)\\
&=-i\epsilon_{ijk}\int d^3x\,\Theta^{0k}
\end{align*}
and therefore
\begin{equation}
[P_j,J_i]=-i\epsilon_{ijk}P_k.
\tag{7.4.19}\label{eq:7.4.19}
\end{equation}
On the other hand, the `boost' generator $K_k\equiv J^{k0}$, though
time-independent, does explicitly involve the time coordinate
\[
K_k=\int d^3x\bigl(x^k\Theta^{00}-x^0\Theta^{0k}\bigr),
\]
or more explicitly
\begin{equation}
\mathbf{K}=-t\mathbf{P}+\int d^3x\,\mathbf{x}\Theta^{00}(\mathbf{x},t).
\tag{7.4.20}\label{eq:7.4.20}
\end{equation}
Since this is a constant, we have
$0=\dot{\mathbf{K}}=-\mathbf{P}+i[H,\mathbf{K}]$, and therefore
\begin{equation}
[H,\mathbf{K}]=-i\mathbf{P}.
\tag{7.4.21}\label{eq:7.4.21}
\end{equation}
Also, applying Eq.~\eqref{eq:7.3.28} again gives
\[
[P_j,K_k]=i\int d^3x\,x^k\frac{\partial}{\partial x^j}\Theta^{00}
=-i\delta_{jk}\int d^3x\,\Theta^{00}
\]
and therefore
\begin{equation}
[P_j,K_k]=-i\delta_{jk}H.
\tag{7.4.22}\label{eq:7.4.22}
\end{equation}
For any reasonable Lagrangian density, the operator \eqref{eq:7.4.20} will be `smooth'
in the sense used in Section 3.3, i.e., the interaction terms in
$e^{iH_0t}\int d^3x\,\mathbf{x}\Theta^{00}(\mathbf{x},0)e^{-iH_0t}$
vanish\footnote{When we say that some interaction-picture operator vanishes
for $t\to\pm\infty$, we mean that its matrix elements between states that are
smooth superpositions of energy eigenstates vanish in this limit.} for
$t\to\pm\infty$. (Note that the interaction terms in
$e^{iH_0t}\int d^3x\,\Theta^{00}(\mathbf{x},0)e^{-iH_0t}$ must vanish for
$t\to\pm\infty$ in order to allow the introduction of `in' and `out' states
and the $S$-matrix.) With this smoothness assumption and the commutation
relation~\eqref{eq:7.4.21} in hand, we can repeat the arguments of Section 3.3, and
conclude that the $S$-matrix is Lorentz-invariant.

\begin{center}
$*\ *\ *$
\end{center}

The same arguments were also used in Section 3.3 to verify that the remaining
commutation relations of the Lorentz group, those of the $J^{\mu\nu}$ with each
other, take the proper form. This can also be shown directly for the
commutators of the rotation generators, which here take the form
\begin{equation}
J^{ij}\equiv\int d^3x\,
\frac{\partial\mathcal{L}}{\partial\dot\Psi^\ell}
\left(x^i\partial_j\Psi^\ell-x^j\partial_i\Psi^\ell
-i(\mathcal{J}^{ij})^\ell{}_m\Psi^m\right).
\tag{7.4.23}\label{eq:7.4.23}
\end{equation}
Since the Lagrangian density does not depend on the time-derivatives of the
auxiliary fields, and the rotation generators do not mix canonical and
auxiliary fields, this can also be written as a sum over canonical fields alone:
\begin{equation}
J^{ij}=\int d^3x\,P_n\left(
x^i\partial_jQ^n-x^j\partial_iQ^n
-i(\mathcal{J}^{ij})^n{}_{n'}\Psi^{n'}\right).
\tag{7.4.24}\label{eq:7.4.24}
\end{equation}
It follows immediately then from the canonical commutation relations that
\begin{equation}
[J^{ij},Q^n(x)]=-i(x_i\partial_j-x_j\partial_i)Q^n(x)
-(\mathcal{J}^{ij})^n{}_{n'}Q^{n'}(x),
\tag{7.4.25}\label{eq:7.4.25}
\end{equation}
\begin{equation}
[J^{ij},P_n(x)]=i(x_i\partial_j-x_j\partial_i)P^n(x)
+(\mathcal{J}^{ij})^{n'}{}_nP_{n'}(x).
\tag{7.4.26}\label{eq:7.4.26}
\end{equation}
These results can be used to derive the usual commutation relations of the
$J^{ij}$ with each other and other generators.\footnote{Also, since $J^{ij}$
commutes with $H$ and $P_n\dot Q^n$, it commutes with $L$. The commutator of
$J^{ij}$ with the auxiliary fields must thus be consistent with the rotational
invariance of $L$.} If there are no auxiliary fields then the same arguments
may be applied to the `boost' generators to complete the demonstration that the
$P^\mu$ and $J^{\mu\nu}$ satisfy the commutation relations of the Poincar\'e
group. However the `boost' matrices $\mathcal{J}^{i0}$ will in general
mix canonical and auxiliary fields (such as the components $V^i$ and $V^0$ of a
vector field), so the direct proof of the commutation relations of the
$J^{i0}$ with each other has to be given on a case by case basis. Fortunately,
this is not needed for the proof of the Lorentz invariance of the $S$-matrix
given in Section 3.3.
